\section{换基：让复杂的作用变成独立的缩放}
\label{sec:17-Synthesis/02-Recurring-Ideas/02}

一份用户表，跑主成分分析。第一主成分几乎完全指向``年收入''这一列，报告里写下``收入是最主要的变异来源''。第二天有人把这一列的单位从元改成万元，重跑同一段代码，第一主成分换成了几乎完全指向``年龄''，前两个主成分张成的平面整个转了方向，散点图看上去像另一份数据。没有人动过算法。

这件事经常被当成 bug 报上来，其实它是定义的直接后果：主成分是样本协方差矩阵的特征向量，而协方差矩阵的元素带量纲，换一个单位就是给矩阵做一次非正交的对角相似变换，特征方向当然会变。真正值得追问的是另一件事——为什么我们会觉得``找一组特征方向''这件事本身有意义，以至于愿意在它给出反常结果时先怀疑代码而不是怀疑做法。

答案是这个动作在这本书里出现得实在太多了。

\subsection{一个动作}

第二个反复出现的数学动作是：\emph{一个纠缠的线性作用，换到合适的坐标里，就变成一组互不干扰的独立缩放。}原来的矩阵 \(A\) 把各个分量搅在一起，换基之后写成 \(A = Q\Lambda Q^{-1}\)，中间那个 \(\Lambda\) 是对角的——在新坐标里，第 \(i\) 个分量只被乘上 \(\lambda_i\)，与其它分量再无关系。于是矩阵的幂、指数、求逆、解方程全部退化成对角线上的标量运算。

\begin{center}
\begin{tikzpicture}[x=1cm,y=1cm,font=\small,>=stealth]
  % ---- 左：原坐标下歪斜的等高线 ----
  \draw[black!35,->] (-2.6,0) -- (2.6,0) node[right,black!55,font=\footnotesize] {\(e_1\)};
  \draw[black!35,->] (0,-1.9) -- (0,1.9) node[above,black!55,font=\footnotesize] {\(e_2\)};
  \begin{scope}[rotate=32]
    \draw[blue!60!black,very thick] (0,0) ellipse (2 and 0.85);
    \draw[blue!30!black,thin] (0,0) ellipse (1.3 and 0.55);
  \end{scope}
  \draw[red!70!black,thick,->] (0,0) -- (1.95,1.22) node[anchor=south west,font=\footnotesize] {\(u_1\)};
  \draw[red!70!black,thick,->] (0,0) -- (-0.69,1.10) node[anchor=south east,font=\footnotesize] {\(u_2\)};
  \node[font=\footnotesize,black!70] at (0,-2.5) {原坐标：两个分量纠缠};

  % ---- 中：换基箭头 ----
  \draw[black!60,->,very thick] (2.95,0) -- (4.15,0);
  \node[font=\footnotesize,black!70,anchor=south] at (3.55,0.12) {\(Q^{\trans}\)};

  % ---- 右：特征基下正的椭圆 ----
  \begin{scope}[shift={(7.0,0)}]
    \draw[red!70!black,->] (-2.6,0) -- (2.6,0) node[right,font=\footnotesize] {\(u_1\)};
    \draw[red!70!black,->] (0,-1.9) -- (0,1.9) node[above,font=\footnotesize] {\(u_2\)};
    \draw[blue!60!black,very thick] (0,0) ellipse (2 and 0.85);
    \draw[blue!30!black,thin] (0,0) ellipse (1.3 and 0.55);
    \draw[black!65,|-|] (0,-1.35) -- (2,-1.35);
    \node[font=\footnotesize,black!75,anchor=north] at (1.0,-1.42) {\(\lambda_1\)};
    \draw[black!65,|-|] (2.3,0) -- (2.3,0.85);
    \node[font=\footnotesize,black!75,anchor=west] at (2.38,0.45) {\(\lambda_2\)};
    \node[font=\footnotesize,black!70] at (0,-2.5) {特征基：两个独立缩放};
  \end{scope}
\end{tikzpicture}

\emph{同一个二次型的同一条等高线。左边看上去是一个需要两个参数描述的歪斜椭圆，右边是两个互不相干的数。换坐标没有改变任何事实，只是把纠缠移进了坐标的选择里。}
\end{center}

\hyperref[sec:03-Linear-Algebra/05-Eigenvalues-and-Eigenvectors/01]{``不变方向与特征值：寻找矩阵最自然的坐标轴''}把这个动作作为特征值概念的出发点：不是先定义 \(Av=\lambda v\) 再问它有什么用，而是先问``有没有一组方向让这个作用变简单''，特征方程是这个问题的形式化。\hyperref[sec:03-Linear-Algebra/05-Eigenvalues-and-Eigenvectors/02]{``对角化与动力系统：把耦合演化变成独立缩放''}给出回报最直接的一处：耦合的线性递推或线性微分方程组，换基后每个坐标独立演化，长期行为完全由谱的位置决定，不必解出显式解就能读出答案。

\subsection{这组好基什么时候存在}

不是每个矩阵都有。\(\begin{pmatrix}1&1\\0&1\end{pmatrix}\) 的特征值只有 \(1\)，特征向量只有 \(e_1\) 的倍数，凑不出一组基，它不可对角化。\hyperref[sec:03-Linear-Algebra/05-Eigenvalues-and-Eigenvectors/03]{``对角化失败与 Jordan 形：独立特征方向不够时怎么办''}说明这种情形下最简形式里必然留下一个上对角线的 \(1\)，而那个 \(1\) 的后果是可见的：\(A^{n}\) 里出现 \(n\lambda^{n-1}\) 这样的多项式因子，纯粹的几何缩放描述不了它。

存在性最干净的一类是对称矩阵。\hyperref[sec:03-Linear-Algebra/06-Symmetric-and-Positive-Definite-Matrices/01]{``对称矩阵与谱定理：什么时候特征方向一定正交''}给出的保证是两层的：特征值全实、特征向量可取成正交基。正交这一层比``存在一组基''值钱得多——正交意味着换基矩阵的逆就是它的转置，换基不放大误差，也不改变长度和角度，于是几何直觉可以原样搬过去。协方差矩阵、Hessian、Laplacian、核矩阵都对称，这不是巧合：它们都由某个双线性的、对两个自变量地位对称的量定义。\hyperref[sec:03-Linear-Algebra/06-Symmetric-and-Positive-Definite-Matrices/04]{``二次曲面用谱分解显露几何形状''}就是上面那张图的完整版本。

一般的矩阵——尤其是不方的矩阵——两条都不满足。\hyperref[sec:03-Linear-Algebra/07-SVD-and-Data-Applications/01]{``任意矩阵的奇异值分解''}给出的出路值得单独指出，因为它常常被当成``对角化的推广''而错过了要点：SVD 之所以对一切矩阵都成立，正是因为它\emph{放弃了输入输出用同一组基}。\(A = U\Sigma V^{\trans}\) 里 \(V\) 是定义域那一侧的正交基，\(U\) 是值域那一侧的，两者一般毫无关系。特征分解要求同一组基既能描述输入又能描述输出，这个额外要求才是它经常做不到的原因。想清楚这一点，``为什么 SVD 总存在而对角化不总存在''就不再是需要记忆的事实。

\subsection{同一个动作的其它名字}

Fourier 变换是这个动作用在一类特殊算子上的结果。所有平移不变的线性算子彼此可交换，因而共享一组特征向量，而那组共同的特征向量就是复指数——这是 Fourier 基的来历，也解释了为什么卷积在频域变成逐点乘法：卷积算子在它自己的特征基下当然是对角的。\hyperref[sec:03-Linear-Algebra/08-Complex-Matrices-and-Fourier-Structure/03]{``Fourier 基与 FFT''}把这条线走完，顺带说明 FFT 是算法而变换本身是数学。反过来读同一件事：\hyperref[sec:14-Deep-Learning/04-Convolutional-Networks/01]{``卷积把平移结构写进线性映射''}里的权值共享，正是在强行要求学到的算子属于这个可以被 Fourier 基对角化的类。

Fourier 基的代价是它在时间上完全离域——每个基函数都铺满整条轴，一次局部的突变会污染所有频率系数。\hyperref[sec:08-Numerical-Analysis/08-Fourier-and-Spectral-Methods/06]{``小波与多尺度分解''}换了一组基，在频率分辨率上让步，换回时间上的局部性。这里出现了这个动作的一般形态：好基不是唯一的，选哪一组取决于你希望哪一类结构变稀疏。

数据分析里的低秩建模是同一个动作加一句额外的断言。\hyperref[sec:13-Machine-Learning/04-Dimensionality-Reduction/02]{``SVD 给出最佳低秩近似''}说明在换到奇异基之后截断前 \(k\) 项，在 Frobenius 与谱范数下都是最优的——注意这句话的力量来自``最优''，而不是来自``可以截断''。同一个断言在别处叫别的名字：微调时的低秩增量假设参数的更新只发生在少数几个方向上；矩阵补全假设观测到的部分足以定出那几个方向。它们共享一个前提，就是那组好基下的系数确实衰减得快，而这个前提是关于数据的经验假设，不是定理。

概率里也有它的影子，只是名字变了。联合分布在一般坐标下需要指数多个参数，而条件独立结构说的正是``存在一种分解，让联合分布写成局部因子的乘积''。\hyperref[sec:13-Machine-Learning/07-Graphical-and-Sequential-Models/02]{``Gaussian 图模型用精度矩阵编码条件独立''}把这个类比坐实到最紧：对 Gaussian 而言，条件独立就是精度矩阵在那一对坐标上的元素为零，图结构就是矩阵的稀疏模式。谱与稀疏在这里是两种不同的``简单''——一个要求换基后对角，一个要求在给定坐标下稀疏，两者一般不能同时达到。\hyperref[sec:13-Machine-Learning/05-Clustering-and-Data-Geometry/02]{``谱聚类把连通性变成特征向量''}走的是前一条路：把图的连通结构塞进 Laplacian 的低端特征向量里，聚类问题于是变成在新坐标下做最朴素的划分。

\subsection{方差最大不是最有用}

开头那个换单位的故事说明了第一件事：主成分方向依赖于量纲，而量纲是人选的。标准化能让结果与单位无关，但标准化本身也是一个选择——它相当于断言每个特征的重要性一样，这个断言未必对。没有``中立''的选项，只有被写明的选项和被隐藏的选项。

第二件事更容易造成实际损失：方差大不等于信息多。判别信息完全可以藏在低方差方向里。两类点沿着一个方向分得清清楚楚但间距很小，同时沿着另一个方向整体拉得很开却两类混在一起，主成分会毫不犹豫地选后者，然后把前者当噪声丢掉。降维之后分类精度反而下降，多数是这么来的。主成分分析是无监督的，它从来没有承诺过保留与标签有关的方向——想保留判别方向就得用有监督的准则，那是另一个优化问题。

第三件事在高维下尤其要紧。样本协方差矩阵的谱不是总体谱的无偏影像：维数与样本量可比时，样本特征值会系统性地向两端散开，最大的偏大、最小的偏小，即便总体协方差是单位阵也会出现一条看上去很像``主要方向''的谱。\hyperref[sec:11-Mathematical-Statistics/09-Random-Matrix-Theory/03]{``谱偏差怎样影响 PCA 与随机投影''}给出这一失真的定量形状。碎石图上的那个拐点，因此可能完全是随机性画出来的。

\subsection{这个动作的边界}

换基不创造信息，它只是重新分配``复杂性''的位置。歪斜的椭圆换到特征基下变成两个数，代价是那两个方向本身要被记下来——省下的和付出的在自由度上完全相等。它带来实际好处的唯一情形，是那组基比数据本身稳定得多：卷积的特征基由平移不变性决定，与数据无关，所以 Fourier 方法可靠；主成分的基由数据估出来，样本一换就变，所以它的结论必须配一句关于稳定性的说明。

这一条与上一篇的判断是同一类。上一篇说线性化只在一个有限邻域内可信，这一篇说好基只在它被稳定确定时可信；两者都是先把问题换成一个便宜的形式，再交代这次换算的有效范围。下一篇处理的情形更彻底：连``换成一个能解的形式''都做不到，只能把解定义成一列近似的极限。
